%% REPLACE SXXXXXX with your student number
\def\studentNumber{S2149988}


%% START of YOUR ANSWERS
%% Add answers to the questions below, by replacing the text inside the brackets {} for \youranswer{ "Text to be replaced with your answer." }. 
%
% Do not delete the commands for adding figures and tables. Instead fill in the missing values with your experiment results, and replace the images with your own respective figures.
%
% You can generally delete the placeholder text, such as for example the text "Question Figure 2 - Replace the images ..." 
%
% There are 18 TEXT QUESTIONS (a few of the short first ones have their answers added to both the Introduction and the Abstract). Replace the text inside the brackets of the command \youranswer with your answer to the question.
%
% There are also 3 "questions" to replace some placeholder FIGURES with your own, and 3 "questions" asking you to fill in the missing entries in the TABLES provided. 
%
% NOTE! that questions are ordered by the order of appearance of their answers in the text, and not by the order you should tackle them. Specifically, you cannot answer Questions 2, 3, and 4 before concluding all of the relevant experiments and analysis. Similarly, you should fill in the TABLES and FIGURES before discussing the results presented there. 
%
% NOTE! If for some reason you do not manage to produce results for some FIGURES and TABLES, then you can get partial marks by discussing your expectations of the results in the relevant TEXT QUESTIONS (for example Question 8 makes use of Table 1 and Figure 2).
%
% Please refer to the coursework specification for more details.


%% - - - - - - - - - - - - TEXT QUESTIONS - - - - - - - - - - - - 



%% Question 1:
\newcommand{\questionOne} {
\youranswer{the training accuracy increases steadily, while validation accuracy reaches a maximum of ~83\% at epoch
12 and then steadily decreases thereafter. We see in figure 1b that the training error decreases throughout,
while validation error reaches a minimum of ~0.5 at epoch 10 and then increases at a steady rate.
Overfitting begins around epoch 12, where the generalisation gap (val. error $-$ train error) increases significiently
at a continued rate. Considering curve velocities, accuracy increases most rapidly during epochs 0 – 8,
then hits a limit very quickly, while validation error’s slope turns from negative to positive at epoch 10.}
}

%% Question 2:
%    Question 2 - Present your network width experiment results by using the relevant figure and table}
\newcommand{\questionTwo} {
\youranswer{the model with 128 hidden units overfits. This can be seen clearly in Figure 2a, as its validation error increases
consistently from epoch 10 onwards. The model, with 64 units, begins to slightly overfit from epochs 25 onwards. Whereas
width 32 doesn't display obvious signs of overfitting, except that it's gneralisation gap increases very slightly over
the epochs.
}
}


%%

%% Question 3:
%    Question 3 - Discuss whether varying width affects the results in a consistent way, and whether the results are expected and match well with the prior knowledge (by which we mean your expectations as are formed from the relevant Theory and literature)}
\newcommand{\questionThree} {
\youranswer{For the parameters tested, increasing the width from 32 to 64 improves performance: the validation accuracy
rises and the validation loss decreases. However, increasing width further to 128 reverses this trend, with both a drop
in validation accuracy and an increase in validation loss. The best validation accuracy (66.6\%) occurs at width 64.
The model that overfits the least is width 32, which has the smallest generalisation gap of 0.154.

This pattern is consistent with expectations from the literature. Increasing width expands model capacity, which
initially helps capture the underlying signal, improving validation accuracy. Beyond a certain point, however, the
additional capacity enables the model to fit noise, increasing overfitting and reducing validation performance.
Width 64 appears to strike the balance in this experiment: large enough to model the structure in the data,
    but not so large that it destabilises generalisation. Width 32, while less accurate, generalises the most reliably.
}
}

%% Question 4:
%    Question 4 - Present your network depth experiment results by using the relevant figure and table}
\newcommand{\questionFour} {
\youranswer{All models show some degree of overfitting, as each displays an increasing generalisation gap across the
epochs. This is most evident in the models with depths 2 and 3. In Figure 3b, their validation error begins to rise from
around epoch 20 onwards, indicating that they start to diverge from the training curve earlier and more noticeably than
the deeper models.
}
}

%% Question 5:
\newcommand{\questionFive} {
\youranswer{For the parameters tested, increasing depth improves validation accuracy but reduces generalisation. Figure 3a shows
that, after 100 epochs, the depth-4 model achieves the highest validation accuracy. Conversely, Figure 3b shows that the
depth-2 model overfits the least, maintaining the smallest generalisation gap across the training run.

These results match expectations from the literature. A network with two hidden layers is already a universal
approximator, provided sufficient width, which explains why depths 2 and 3 achieve strong baseline performance.
However, additional depth increases model capacity further, enabling the network to capture more complex structure—but
also making it more prone to fitting noise. This aligns with the bias–variance trade-off: deeper networks can model
richer functions but at the cost of greater overfitting unless regularised appropriately.
}
}


%% Question 6:
%    Question 6 - Explain how one could use a combination of L1 and L2 regularisation. Discuss any potential benefits of this approach. Present an equation for this combined loss function and explain all relevant variables and hyperparameters.}
\newcommand{\questionSix} {
\youranswer{We combine the L1 and L2 penalties together by adding them both to the loss function. By supplying a
coefficient alpha, we can balance between the sparsity of L1 and the smoothing of L2. This can be implemented by
wrapping our standard loss function like the equation below:

\[
L_{\text{total}}
= L_{\text{data}}
+ \lambda \left[
    \alpha \, \|w\|_{1}
    + (1 - \alpha)\, \|w\|_{2}^{2}
\right].
\]

\begin{itemize}
    \item \(L_{\text{total}}\): The complete loss the optimiser seeks to minimise.
    \item \(L_{\text{data}}\): The data-driven loss term (e.g.\ cross-entropy or MSE).
    \item \(w\): The vector of model parameters being regularised.
    \item \(\|w\|_{1}\): L1 penalty, encouraging sparsity by pushing weights towards zero.
    \item \(\|w\|_{2}^{2}\): L2 penalty, encouraging smoother and smaller weights.
    \item \(\lambda\): The overall regularisation strength that scales the combined penalty.
    \item \(\alpha \in [0,1]\): Balancing coefficient controlling the mix of L1 and L2.
    \item \(1 - \alpha\): Complementary weight ensuring the two penalties sum cleanly.
\end{itemize}

The additional hyperparameter \(\alpha\), is important because it directly controls the trade-off between sparsity and
smoothness. The literature supports this reasoning: Ng~\cite{ng2004feature} shows that L1 is more
aggressive than L2 in driving weights to zero, especially when features are correlated, while Goodfellow et
al.~\cite{Goodfellow-et-al-2016} discuss how L2 regularisation stabilises optimisation by discouraging large weights.
Combining L1 and L2 therefore allows us to retain the advantages of both. L1 can remove genuinely irrelevant weights,
while L2 prevents overly aggressive elimination of correlated but still informative parameters.

This mixed approach could provide a more flexible and robust regulariser than using L1 or L2 alone, potentially leading
to better generalisation performance depending on the structure of the data.
}
}



%% Question 7:
%    Question 7 - Assume you had a limited experiment budget for exactly \textbf{4} experiment runs dedicated to testing the effectiveness of your proposed combination of L1 and L2 Regularization in Section~\ref{sec:L1L2combo}. Argue for which 4 configurations to run based on any previous results or other information you consider relevant. Run those experiments and add the results in Table~4 and Figure~5 (you will then discuss those results alongside all others in the next question below).
\newcommand{\questionSeven} {
\youranswer{
Due to limited experimental budget of four runs, we focused on the interaction between L1 and L2 rather than
just finding a single “best” setting. From our earlier tests, both L1 and L2 have useful ranges around 5e-4-1e-3,
with performance peaking near 1e-11 and degrading when too large.
Guided by that, we probed two balanced points and two asymmetric points to test whether gains are driven primarily by one term or if a complementary effect emerges.

    Specifically, we ran:
    \begin{itemize}
        \item Balanced small: \(\lambda_1=5\text{e-}4,\ \lambda_2=5\text{e-}4\) — tests a minimal joint penalty near each method’s individual sweet spot
        \item Balanced moderate: \(\lambda_1=1\text{e-}09,\ \lambda_2=5\text{e-}3\) — tests a stronger joint penalty to see if combined effects scale
        \item Stronger L1: \(\lambda_1=5\text{e-}3,\ \lambda_2=5\text{e-}4\) — tests the effect of a dominant L1 term
        \item Stronger L2: \(\lambda_1=5\text{e-}4,\ \lambda_2=5\text{e-}3\) — tests the effect of a dominant L2 term
    \end{itemize}

    This set stays within the four-run budget while answering a clearer question: do we see synergy between L1 and L2,
    or does one term dominate? Using balanced vs asymmetric pairs lets us attribute any improvement to either
    interaction effects or to a single regulariser, and the chosen magnitudes reflect the effective region indicated by
    our prior results.
}
}


%% Question 8:
\newcommand{\questionEight} {
\youranswer{
Across all experiments, each regularisation method reduces overfitting to some extent, but only for suitable
hyperparameter values. Once these values become too large, performance degrades sharply. From our results, Dropout
with an inclusion rate of $0.97$ achieves the highest validation accuracy at $84.9\%$. The least overfitting model in
absolute terms is L1 regularisation with parameters $5\times10^{-2}$ and $5\times10^{-3}$, both producing a
generalisation gap of $0$. However, this is misleading because the corresponding validation accuracy is only $2\%$,
which is effectively random chance for a 47-class problem ($1/47 \approx 0.0213$). Excluding all models with validation
accuracy below $70\%$, the smallest overfitting still occurs under L1 regularisation at $\lambda = 5\times10^{-4}$,
with a generalisation gap of $0.014$ ($0.824 - 0.810$). This suggests that very mild sparsity may remove unhelpful
weights without harming useful structure.

Comparing all methods to the baselines, regularisation improves validation accuracy and reduces error whenever
reasonable hyperparameter values are chosen. The main exception is L1, which fails to outperform the baseline at any of
the tested values. Soft labelling ($\tau = 0.1$) lowers validation accuracy but improves the generalisation gap
($0.156$ vs.\ the baseline’s $0.242$), indicating reduced overfitting at the expense of predictive performance.

Comparing L1 and L2, L2 is consistently superior for $\lambda > 10^{-3}$, achieving higher validation accuracy and
lower error rates. This aligns with established findings that L2 regularisation yields smoother optimisation behaviour
and discourages extreme weight growth \cite{Goodfellow-et-al-2016,ng2004feature}. However, the smallest L1 values still
produce the best generalisation gap among the single-regulariser experiments while retaining reasonable accuracy.

Overall, overfitting is reduced by all methods tested, so long as the hyperparameters are not excessively large. The
best-performing model in terms of accuracy is Dropout at $p = 0.85$, and this is further supported by the observation
that the next-best method is L2, which smooths weights and distributes reliance across hidden units. This connection is
consistent with the interpretation of dropout as a stochastic mechanism that limits co-adaptation and regulates weight
growth in a way similar to L2 \cite{srivastava2014dropout}.

The combined L1--L2 experiments produce two further conclusions. First, balancing L1 and L2 yields both low
generalisation gaps and moderate-to-high validation accuracies. Second, favouring L1 within the combination increases
validation accuracy but also increases error, indicating a trade-off between sparsity and stability. Notably, the
best generalisation gap across all models ($0.006$) is achieved by an L1--L2 combination, suggesting that mixed
regularisation is the most effective method for reducing overfitting while maintaining acceptable accuracy. However,
its accuracy remains below the best L2 and Dropout settings, showing that reducing overfitting does not always
translate into the highest predictive performance.

Finally, the best-performing model (Dropout at $p = 0.85$) was trained and evaluated on the test set, achieving a test
accuracy of $0.854$ and test error of $0.492$ (reported in Table~3). This confirms that the improvements observed during
validation also translate to unseen data.
}
}


%% Question 9:
%    Question 9 - A colleague proposed using a custom activation function (which can be identified as CustomActivationLayer() in your Coursework\_1.ipynb notebook and the mlp package). Run the prepared experiment in the "Problematic Training" cell in the Coursework\_1.ipynb notebook, and display the results in Figure~6. Analyse the results. Is the model overfitting? Explain the cause of the problematic performance of the model?
\newcommand{\questionNine} {
\youranswer{
A colleague proposed training a depth-5 model using their custom activation function
\(
f(x) = \frac{1}{10(1 + e^{-x})}.
\)
Our experiment shows that the model is severely underfitting. Figure~6a and Figure~6b clearly show that the
validation accuracy remains low and the validation error remains high throughout training, with no indication of the
model learning meaningful structure. Figure~6c further confirms the cause: the gradients collapse in the upper layers of
the network during backpropagation.

This problematic behaviour is directly caused by the activation function itself. The proposed function is effectively a
scaled sigmoid with outputs restricted to the extremely narrow range $(0, 0.1)$. These heavy compression forces
activations to cluster tightly, producing gradients that are already very small even before backpropagation begins.
During training, these gradients diminish further, resulting in a classic vanishing-gradient failure mode. Because the
updates to the weights are negligible, the network cannot fit the data, leading to persistent underfitting.

To verify this, we reran the experiment using a standard ReLU activation. In this case, gradients remained non-zero
across the network and the model successfully learned. This demonstrates that the poor performance arises from the
activation choice rather than the model depth or data. The issue can be resolved by replacing the custom activation with
one that preserves gradient magnitude, such as ReLU or another non-saturating function.
}

}




%% Question 10:
\newcommand{\questionTen} {
\youranswer{
Across all experiments, the results show that model capacity and regularisation are the primary factors determining
generalisation. Increasing width or depth initially improves validation accuracy, but once the network becomes too
large, the generalisation gap widens. This matches the expected capacity--overfitting behaviour outlined in
\cite{Goodfellow-et-al-2016}, where larger networks can represent more complex structure but also fit noise more
readily.

Across the regularisation methods tested, Dropout was the most consistently effective. An inclusion rate of $0.85$
achieved the highest validation accuracy, which aligns with \cite{srivastava2014dropout}, where dropout is shown to
reduce co-adaptation and encourage more robust representations. L2 regularisation also improved accuracy and reduced
overfitting, consistent with the stabilising effect on weight magnitudes described in
\cite{Goodfellow-et-al-2016,ng2004feature}. L1 regularisation only helped when the coefficient was very small; as
\cite{ng2004feature} explains, larger L1 penalties aggressively zero correlated or moderately informative weights, which
in our results led to severe underfitting. The combined L1--L2 regulariser produced the smallest generalisation gap in
the entire study ($0.006$), demonstrating that mixing sparsity and smoothness can be highly effective for controlling
overfitting even if it does not yield the top validation accuracy.

Comparing Dropout-only with Weight-only approaches, the experimental results show that Dropout is superior for
validation accuracy, whereas L2 and mild L1 excel at reducing the generalisation gap. This indicates that overfitting in
our task arises from both excessive dependence on individual hidden units (addressed by Dropout) and uncontrolled weight
growth (addressed by L2). Together, these results demonstrate that the different regularisers act on different failure
modes, and that combining them can provide complementary benefits.

Beyond the techniques tested, several additional approaches could be explored. Data augmentation and early stopping,
both recommended in \cite{Goodfellow-et-al-2016}, are standard and effective mechanisms for reducing overfitting by
either increasing the effective dataset size or halting training before memorisation occurs. These would provide
natural extensions to the present study.

Our strongest validation accuracy was achieved using Dropout, while the smallest generalisation gap came from the
combined L1--L2 method. This raises an interesting future direction: whether we can achieve the strong accuracy of
Dropout while maintaining the stability seen with L1--L2. One promising option is to investigate activation functions
designed to pair well with Dropout, such as maxout \cite{goodfellow2013maxout}, which preserves gradient flow and may
support higher accuracy without sacrificing generalisation. Another practical extension is to adopt stronger data
augmentation pipelines or incorporate early stopping to further limit overfitting without altering the architecture.

Overall, these findings reinforce that network capacity must be controlled deliberately. Both width and depth expand the
hypothesis space, increasing the risk of overfitting, while the choice of regulariser determines how effectively this
capacity is constrained. Our experiments show that Dropout, L2, and combined L1--L2 each target different dimensions of
overfitting, emphasising the importance of matching the regularisation strategy to the architecture and the desired
balance between accuracy and generalisation.
}
}

%% - - - - - - - - - - - - FIGURES - - - - - - - - - - - - 

%% Question Figure 2:
%    Question Figure 2 - Replace the images in Figure 2 with figures depicting the accuracy and error, training and validation curves for your experiments varying the number of hidden units.
\newcommand{\questionFigureTwo} {
\youranswer{
\begin{figure}[t]
    \centering
    \begin{subfigure}{\linewidth}
        \includegraphics[width=\linewidth]{figures/width_acc}
        \caption{accuracy by epoch}
        \label{fig:width_acccurves}
    \end{subfigure} 
    \begin{subfigure}{\linewidth}
        \centering
        \includegraphics[width=\linewidth]{figures/width_error}
        \caption{error by epoch}
        \label{fig:width_errorcurves}
    \end{subfigure} 
    \caption{Training and validation curves in terms of classification accuracy (a) and cross-entropy error (b) on the EMNIST dataset for different network widths.}
    \label{fig:width}
\end{figure} 
}
}

%% Question Figure 3:
%    Question Figure 3 - Replace these images with figures depicting the accuracy and error, training and validation curves for your experiments varying the number of hidden layers.
%
\newcommand{\questionFigureThree} {
\youranswer{}
\begin{figure}[t]
    \centering
    \begin{subfigure}{\linewidth}
        \includegraphics[width=\linewidth]{figures/depth_acc}
        \caption{accuracy by epoch}
        \label{fig:depth_acccurves}
    \end{subfigure} 
    \begin{subfigure}{\linewidth}
        \centering
        \includegraphics[width=\linewidth]{figures/depth_error}
        \caption{error by epoch}
        \label{fig:depth_errorcurves}
    \end{subfigure} 
    \caption{Training and validation curves in terms of classification accuracy (a) and cross-entropy error (b) on the EMNIST dataset for different network depths.}
    \label{fig:depth}
\end{figure} 
}

%% Question Figure 4:
%    Question Figure 4 - Replace these images with figures depicting the Validation Accuracy and Generalisation Gap (difference between validation and training error) for each of the experiment results varying the Dropout inclusion rate, and L1/L2 weight penalty depicted in Table 3 (including any results you have filled in).
\newcommand{\questionFigureFour} {
\youranswer{
\begin{figure*}[t]
    \centering
    \begin{subfigure}{.475\linewidth}
        \includegraphics[width=\linewidth]{figures/dropout_plot}
        \caption{Accuracy and error by inclusion probability.}
        \label{fig:dropoutrates}
    \end{subfigure} 
    \begin{subfigure}{.475\linewidth}
        \centering
        \includegraphics[width=\linewidth]{figures/l2_plot}
        \caption{Accuracy and error by weight penalty.}
        \label{fig:weightrates}
    \end{subfigure} 
    \caption{Accuracy and error by regularisation strength of each method (Dropout and L1/L2 Regularisation).}
    \label{fig:hp_search}
\end{figure*}
}
}

%% Question Figure 5:
%    Question Figure 5 - Add figures depicting the Validation Accuracy and Generalisation Gap (difference between validation and training error) for each of your 4 experiments  varying your hyperparameters for the combination of L1 and L2 regularisation (based on your implementation) depicted in Table 4.
\newcommand{\questionFigureFive} {
\youranswer{
\begin{figure*}[t]
    \centering
    \begin{subfigure}{.475\linewidth}
        \includegraphics[width=\linewidth]{figures/q5_all_l1l2_valacc}
        \caption{Validation Accuracy Across L1L2 Configurations}
        \label{fig:L1L2valacc}
    \end{subfigure}
    \begin{subfigure}{.475\linewidth}
        \includegraphics[width=\linewidth]{figures/q5_all_l1l2_gengap}
        \caption{Generalisation Gap Across L1L2 Configurations}
        \label{fig:L1L2gengap}
    \end{subfigure}
\end{figure*}
}
}

%% Question Figure 6:
%    Question Figure 6 - Add figures depicting the Training Accuracy, Error, and Gradient Flow Across Layers for the 1 experiment run with the Custom Activation function (produce these by running the "Problematic Training" cell in the Coursework\_1.ipynb notebook).
\newcommand{\questionFigureSix} {
\youranswer{
\begin{figure}[t]
    \centering
    \begin{subfigure}{\linewidth}
        \includegraphics[width=\linewidth]{figures/custom_activation_acc}
        \caption{accuracy by epoch}
        \label{fig:custom_act_acccurves}
    \end{subfigure}
    \begin{subfigure}{\linewidth}
        \centering
        \includegraphics[width=\linewidth]{figures/custom_activation_error}
        \caption{error by epoch}
        \label{fig:custom_act_errorcurves}
    \end{subfigure}
    \begin{subfigure}{\linewidth}
        \centering
        \includegraphics[width=\linewidth]{figures/custom_activation_grad_flow}
        \caption{gradient flow across layers}
        \label{fig:custom_act_gradflow}
    \end{subfigure}
    \caption{Training accuracy, error and gradient flow for a custom activation function}
    \label{fig:custom_act}
\end{figure}

}
}

%% - - - - - - - - - - - - TABLES - - - - - - - - - - - - 

%% Question Table 1:
\newcommand{\questionTableOne} {
\youranswer{
\begin{table}[t]
    \centering
    \begin{tabular}{c|ccc}
    \toprule
        \# Hidden Units & Val. Acc. & Train Error & Val. Error\\
    \midrule
    32   &   0.583   &   1.633   &   1.787   \\ %0.154
    64   &   0.666   &   1.087   &   1.365   \\ %0.278
    128   &   0.604   &   1.172   &   1.721   \\ %0.549
    \bottomrule
    \end{tabular}
    \caption{Validation accuracy (\%) and training/validation error (in terms of cross-entropy error) for varying network widths on the EMNIST dataset.}
    \label{tab:width_exp}
\end{table}
}
}

%% Question Table 2:
%Question Table 2 - Fill in Table 2 with the results from your experiments varying the number of hidden layers.
%
\newcommand{\questionTableTwo} {
\youranswer{
\begin{table}[t]
    \centering
    \begin{tabular}{c|ccc}
    \toprule
        \# Hidden Layers & Val. Acc. & Train Error & Val. Error \\
    \midrule
        1   &   0.537   &   1.771   &   1.904   \\
        2   &   0.584   &   1.164   &   1.762   \\
        3   &   0.641   &   1.104   &   2.026   \\
    \bottomrule
    \end{tabular}
    \caption{Validation accuracy (\%) and training/validation error (in terms of cross-entropy error) for varying network depths on the EMNIST dataset.}
    \label{tab:depth_exps}
\end{table}
}
}

%% Question Table 3:
%Question Table 3 - Fill in Table 3 with the results from your experiments for the missing hyperparameter values for each of L1 regularisation, L2 regularisation, Dropout and label smoothing (use the values shown on the table).
\newcommand{\questionTableThree} {
\youranswer{
\begin{table*}[t]
    \centering
    \begin{tabular}{c|c|ccc}
    \toprule
        Model    &  Hyperparameter value(s) & Validation accuracy & Train Error & Validation Error \\
    \midrule
    \midrule
        Baseline &  -         &   0.763   &   0.676   &   0.918   \\
    \midrule
        \multirow{4}*{Dropout}
                &    0.6   &   0.799   &   1.472   &   1.507   \\
                &    0.7   &   0.832   &   0.902   &   0.947   \\
                &    \textbf{0.85}   &  \textbf{0.851}  &   \textbf{0.414}   &   \textbf{0.497}   \\
                &    0.97   &   0.849   &   0.332   &   0.509   \\
    \midrule
        \multirow{4}*{L1 penalty}% 5e-4, 1e-3, 5e-3, 5e-2
                &    \emph{5e-4}   &   \emph{0.763}   &   \emph{0.810}   &   \emph{0.824}   \\
                &    1e-3   &   0.701   &   1.064   &   1.084   \\
                &    5e-3   &   0.020   &   3.850   &   3.850   \\
                &    5e-2   &   0.020   &   3.850   &   3.850   \\
    \midrule
        \multirow{4}*{L2 penalty}  
                &    5e-4   &   0.748   &   0.823   &   0.949   \\
                &    \emph{1e-3}   &   \emph{0.842}   &   \emph{0.451}   &   \emph{0.539}   \\
                &    5e-3   &   0.772   &   0.833   &   0.852   \\
                &    5e-2   &   0.332   &   2.567   &   2.568   \\
    \midrule
            Label smoothing & 0.1 & 0.747   &   1.352   &   1.508   \\
    \bottomrule
    \end{tabular}
    \caption{Results of all hyperparameter search experiments. \emph{italics} indicate the best results per series (Dropout, L1 Regularisation, L2 Regularisation, Label smoothing) and \textbf{bold} indicates the best overall.}
    \label{tab:hp_search}
\end{table*}
}
}


%% Question Table 4:
\newcommand{\questionTableFour} {
\youranswer{
\begin{table*}[t]
    \centering
    \begin{tabular}{cc|ccc}
    \toprule
    \ L1 Coefficient & L2 Coefficient & Val. Acc. & Train Error & Val. Error \\
    \midrule
    5e-4   &   5e-4   &   0.772   &   0.775   &   0.781   \\ % gen gap = 0.006
    5e-4   &   1e-3   &   0.759   &   0.883   &   0.889   \\ % gen gap = 0.006
    1e-3   &   5e-4   &   0.690   &   1.135   &   1.154   \\ % gen gap = 0.019
    1e-3   &   1e-3   &   0.664   &   1.287   &   1.305   \\ % gen gap = 0.018
    \bottomrule
    \end{tabular}
\end{table*}
}
}

%% END of YOUR ANSWERS
